\documentclass[11pt]{article}
\usepackage[a4paper,margin=28mm]{geometry}
\usepackage[T1]{fontenc}
\usepackage{lmodern,amsmath,amssymb,amsthm,mathtools,hyperref}
\newcommand{\ddc}{dd^c}
\newcommand{\dc}{d^c}
\newcommand{\MA}{\operatorname{MA}}
\newcommand{\PSH}{\operatorname{PSH}}
\newcommand{\Eenv}{\mathcal E_{\rm env}}
\newtheorem{theorem}{Theorem}
\newtheorem{lemma}[theorem]{Lemma}
\theoremstyle{remark}\newtheorem{remark}[theorem]{Remark}
\title{A logarithmic-potential criterion for singular measures at zero lower mass}
\author{}\date{}
\begin{document}\maketitle
\noindent\textbf{Working draft.} The measure trace, variational argument,
primary applications and product conclusion have not passed native
review, independent reconstruction and final comparison.
\begin{abstract}
On a product with a Hermitian factor of zero lower Monge--Amp\`ere mass,
we give a sufficient condition on an individual measure using the
explicit logarithmic kernel of the projective-line factor. Every non-polar
curve probability supplies an actual normalized finite-mass equation;
boundedness and continuity have exact linear-potential tests in the
separated sector. Under the finite continuous-potential condition, a
convex Hilbert-space functional with a proved canonical measure trace
supplies every specified positive-exponent equation.
The entire centered curve family has a uniform selected limit with an
order characterization. The condition includes nonconstant weighted
arclength probabilities, singular with respect to ambient smooth volume,
and positive-dimensional Frostman measures. The three equations retain
their different normalizations. The product coefficient is maximal over all
genuine subsolutions, and the bounded exponential solution is greatest even
against singular competitors. No necessary criterion for arbitrary measures
is asserted.
\end{abstract}

Use $\ddc=i\partial\bar\partial/(2\pi)$,
$\dc=i(\bar\partial-\partial)/(4\pi)$ and ordinary wedge powers.
Potentials always have their genuine canonical values. Extended products
are non-pluripolar products with smooth-background corrections.

\section{The criterion and the three equations}
Let $(X_0,h_0)$ be a connected compact Hermitian $m$-fold with
$v_-(h_0)=0$, and put $V_0=\int_{X_0}h_0^m$.
On $X=X_0\times Y$, $Y=\mathbf P^1$, put
\[
 N=m+1,\qquad H=h_0+\kappa,\qquad \int_Y\kappa=1,
\]
where $\kappa$ is Fubini--Study and pullbacks are implicit.
For the chordal distance define
\[
 G(x,y)=\log d_{\rm ch}(x,y)^2+1,\qquad K=1-G\geq0.
\]
Its normalization is
$\ddc_xG(x,y)=\delta_y-\kappa$ and $\int_YG(x,y)\kappa(x)=0$.
Let $\mu$ be a Borel probability on $Y$ satisfying the explicit criterion
\begin{equation}\label{eq:criterion}
 Q_\mu(x)=\int_YG(x,y)\,d\mu(y)\quad\hbox{is finite and continuous on }Y.
\end{equation}
For any $M>0$ prescribe the geometric product measure
\[
 \nu=M\frac{h_0^m}{V_0}\otimes\mu,\qquad c=\frac{NV_0}{M}.
\]
Condition \eqref{eq:criterion} tests a linear kernel integral of the
measure, without assuming existence of a nonlinear subsolution.

\begin{theorem}\label{thm:main}
The measure $\nu$ is finite, nonzero and non-pluripolar. The background
$H$ is nonclosed and has $v_-(H)=0$. The pullback of
\[
 u=Q_\mu-\sup_YQ_\mu
\]
is bounded continuous, genuine $H$-psh, and satisfies
\[
 \sup_Xu=0,\qquad \MA_H^{\rm ext}(u)=c\nu,\qquad
 \int_X\MA_H^{\rm ext}(u)=NV_0.
\]
In particular every $0<a\leq c$ has a positive scaled subsolution.

For each specified $\lambda>0$ there is a unique bounded genuine curve
solution $w_\lambda$, which is continuous and satisfies
\[
 \kappa+\ddc w_\lambda=e^{\lambda w_\lambda}\mu,\qquad
 \int_Ye^{\lambda w_\lambda}\,d\mu=1.
\]
Then $U_\lambda=\lambda^{-1}\log c+w_\lambda\circ\pi_Y$ solves
\[
 \MA_H^{\rm ext}(U_\lambda)=e^{\lambda U_\lambda}\nu,\qquad
 \sup_XU_\lambda=\lambda^{-1}\log c+\sup_Yw_\lambda.
\]
Its supremum-zero normalization $v_\lambda$ therefore has the different
equation
\[
 \MA_H^{\rm ext}(v_\lambda)
   =c e^{\lambda\sup w_\lambda}e^{\lambda v_\lambda}\nu.
\]
All these bounded potentials belong to $\Eenv(X,H,0)$.
The solution $U_\lambda$ is the greatest genuine exponential subsolution,
including competitors with extra poles, excess mass or infinite mass.
It is the unique bounded product solution. The coefficient $c$ is the
attained maximum of $a\geq0$ with $\MA_H^{\rm ext}(p)\geq a\nu$ over
all genuine $H$-psh potentials. The same maximum is attained in the bounded
ordinary type and in $\Eenv(X,H,0)$.

As $\lambda\downarrow0$, the entire family satisfies
\[
 w_\lambda\longrightarrow Q_\mu-\int_YQ_\mu\,d\mu,
 \qquad v_\lambda\longrightarrow u
\]
uniformly, with error $O(\lambda)$. Within the separated sector the
first limit is the greatest critical subsolution of $\mu$-mean at most
zero; $u$ is the unique supremum-zero critical member. This asserts no
uniqueness for arbitrary nonclosed critical potentials or unbounded
positive-exponent solutions.
\end{theorem}

\section{Finite equations without the continuity condition}
\begin{lemma}\label{lem:arbitrary}
Retain the product geometry and normalization, but replace
\eqref{eq:criterion} by the assumption that $\mu$ is any probability
giving no mass to polar Borel sets on $Y$. Then $Q_\mu$ is a genuine
$\kappa$-psh function of mean zero, possibly unbounded, and
$u=(Q_\mu-\sup_YQ_\mu)\circ\pi_Y$ solves the same exact finite equation
$\MA_H^{\rm ext}(u)=c\nu$. It is the unique supremum-zero solution in
the sector of pullbacks of genuine $\kappa$-psh functions.
Such a separated solution is bounded if and only if $Q_\mu$ is bounded,
and is continuous if and only if $Q_\mu$ is continuous.
\end{lemma}
\begin{proof}
Set $G_L=\max(G,-L)$. This is a bounded continuous kernel and, for
each fixed $y$, a genuine $\kappa$-psh function in $x$. The integrals
$Q_L=\int G_L\,d\mu$ are continuous $\kappa$-psh functions decreasing
to the displayed $Q_\mu$. Since $G\leq1$ and
$\int_Y|G(x,y)|\kappa(x)$ is uniformly finite, Tonelli shows that
$Q_\mu$ is integrable, has mean zero and is not identically $-\infty$.
The decreasing local psh theorem retains its canonical values, and
distributional Fubini gives $\kappa+\ddc Q_\mu=\mu$.

For a genuine curve potential $p$, the degree-one non-pluripolar
product is $1_{\{p>-\infty\}}(\kappa+\ddc p)$
\cite[Section 1.2]{BEGZ}. The pole set of $Q_\mu$ is polar, hence has
$\mu$-mass zero. Its actual product is therefore $\mu$.
For $p_j=\max(p,-j)$, the bounded corrected-product expansion gives
\[
 \MA_H(p_j\circ\pi_Y)=N h_0^m\otimes(\kappa+\ddc p_j).
\]
Restrict to $\{p>-j\}$ and increase $j$. The genuine finite-level
definition proves the same formula with the degree-one non-pluripolar
product of $p$ on the right. Applying this to $Q_\mu$ proves the
claimed equation and mass $NV_0$. Non-pluripolarity of $\nu$ follows
from curve slicing and two applications of Fubini, as below; it uses
only the assumed non-polarity of $\mu$ here.

Conversely, a separated solution has degree-one non-pluripolar
measure $\mu$ of mass one. Its entire positive curve current also
has mass one by distributional Stokes. Thus no positive polar part
remains, and $\kappa+\ddc p=\mu$. Harmonic distributional uniqueness
and canonical representatives imply $p=Q_\mu+C$. This proves all
three characterization assertions. No boundedness, stretching
membership or exponential existence is inferred for a general
unbounded $Q_\mu$ from this lemma.
\end{proof}

\section{Non-pluripolarity and the canonical measure trace}
Write $Q=Q_\mu$. The criterion gives
\begin{equation}\label{eq:kernelbound}
 \sup_x\int_YK(x,y)\,d\mu(y)=\sup_x(1-Q(x))<\infty.
\end{equation}
On smaller plane coordinate disks this kernel differs from
$-\log|z-a|^2$ by a uniformly bounded smooth term. The local Riesz
decomposition and canonical uniqueness in \cite[I.4.20--4.21]{Dem}
write a nontrivial subharmonic function as a harmonic remainder plus
the logarithmic potential of a compactly supported finite positive
measure, near a smaller compact disk. Tonelli and
\eqref{eq:kernelbound} make that potential integrable against $\mu$;
the harmonic remainder is bounded there. Thus $\mu$ gives zero mass
to every local polar Borel set.

In a coordinate product, a nontrivial local psh function has nontrivial
subharmonic curve restrictions for smooth-volume-almost every point of
$X_0$, by local integrability and Fubini. The preceding fact and a
second Fubini application show that $\nu$ gives no mass to its pole
set. Countably many smaller product charts prove non-pluripolarity.
No ambient volume density or inference from zero Lelong numbers is used.

For real functions define the Dirichlet form
\[
 E(p,q)=\int_Y dp\wedge\dc q.
\]
Smooth rotation convolution on $\mathrm{SU}(2)$ preserves $\kappa$ and
gives $Q_\epsilon\to Q$ uniformly, probabilities $\mu_\epsilon$ and
$\ddc Q_\epsilon=\mu_\epsilon-\kappa$, with $\kappa$-mean zero.
Smooth Stokes gives
\[
 E(Q_\epsilon,Q_\epsilon)
 =-\int_YQ_\epsilon\,d\mu_\epsilon\leq\|Q\|_\infty.
\]
Poincar\'e and weak Hilbert compactness imply $Q\in H^1(Y)$, and
distributional passage to the limit gives, for smooth real $h$,
\begin{equation}\label{eq:traceidentity}
 \int_Yh\,d\mu=\int_Yh\kappa-E(h,Q).
\end{equation}
Apply this to $\sqrt{h^2+\epsilon^2}$ and use its gradient bound by
$|dh|$. Passing to zero gives
\begin{equation}\label{eq:tracebound}
 \int_Y|h|\,d\mu\leq C_\mu\|h\|_{H^1}.
\end{equation}
Completion therefore defines a continuous trace
$H^1(Y)\to L^1(\mu)$.

For a genuine $\kappa$-psh $H^1$ function this trace equals its canonical
restriction. Indeed take rotation kernels invariant under point
stabilizers, so the corresponding small-ball averages are radial.
They converge strongly in $H^1$ and pointwise to the canonical function:
locally add a smooth quadratic psh correction, use the submean inequality
for the lower limit, and upper semicontinuity for the upper limit.
The correction's averaging error tends to zero. This works also at
minus-infinity points. Bound \eqref{eq:tracebound} gives $L^1(\mu)$
convergence, and a subsequence has $\mu$-a.e. convergence to the trace.
Non-polarity makes the canonical function finite $\mu$-a.e., identifying
the two. Arbitrary changes on volume-null sets are excluded.

\section{Existence for every specified positive exponent}
Fix $\lambda>0$. On the real Hilbert space
\[
 \mathcal H=\{q\in H^1(Y):\int_Yq\kappa=0\}
\]
consider the extended functional, using the trace above,
\[
 J_\lambda(q)=\frac12E(q,q)
             +\frac1\lambda\log\int_Ye^{\lambda q}\,d\mu.
\]
It is finite at zero. Jensen and \eqref{eq:traceidentity} give
\[
 J_\lambda(q)\geq\frac12E(q,q)+\int_Yq\,d\mu
                =\frac12E(q,q)-E(q,Q).
\]
Thus it is coercive and bounded below. It is strongly lower
semicontinuous: strong $H^1$ convergence gives $L^1(\mu)$ convergence;
a subsequence realizing the liminf converges $\mu$-a.e., and Fatou
applies to its positive exponentials. Jensen excludes zero integrals
and minus-infinite logarithms. H\"older proves convexity of the
logarithmic term. The Dirichlet term is strictly convex on
$\mathcal H$. A closed convex epigraph is weakly closed by separation;
weak Hilbert compactness therefore gives a unique minimizer $q_\lambda$
with $0<Z_\lambda=\int e^{\lambda q_\lambda}\,d\mu<\infty$.

For each smooth mean-zero $h$, differentiate at $q_\lambda$ in direction
$h$. Boundedness of $h$ supplies an integrable domination by a constant
times $e^{\lambda q_\lambda}$. The variational equation is
\[
 E(q_\lambda,h)+\int_Yh\frac{e^{\lambda q_\lambda}}{Z_\lambda}\,d\mu=0.
\]
Writing $\rho_\lambda=e^{\lambda q_\lambda}/Z_\lambda$, this is
\begin{equation}\label{eq:Euler}
 \ddc q_\lambda=\rho_\lambda\mu-\kappa,\qquad
 \int_Y\rho_\lambda\,d\mu=1.
\end{equation}
The constants are fixed by total masses. It follows that $q_\lambda$
has a genuine canonical $\kappa$-psh representative, and its trace is
that representative $\mu$-a.e. Its mean-zero Green representation is
\[
 q_\lambda=G*(\rho_\lambda\mu),\qquad q_\lambda\leq1.
\]
Distributional equality gives pointwise equality of genuine
representatives. Hence
$\rho_\lambda\leq e^\lambda/Z_\lambda$ $\mu$-a.e. Bound
\eqref{eq:kernelbound} gives a finite lower bound for $q_\lambda$;
it is bounded.

It is continuous as well. The continuous functions
$x\mapsto\int\min(K(x,y),L)\,d\mu(y)$ increase to the continuous
function $1-Q(x)$. Dini's theorem gives
\begin{equation}\label{eq:tail}
 \sup_x\int_Y(K(x,y)-L)_+\,d\mu(y)\longrightarrow0.
\end{equation}
The same estimate holds for the measure $\rho_\lambda\mu$, because
$\rho_\lambda$ is bounded. Its continuous truncated-kernel potentials
converge uniformly, proving continuity of $q_\lambda$.

Set $w_\lambda=q_\lambda-\lambda^{-1}\log Z_\lambda$.
Equation \eqref{eq:Euler} is the scalar-one curve equation in the
theorem, including its additive constant. Any bounded genuine curve
solution is in $H^1$, by the same smooth-energy argument, and its
trace is canonical. For two bounded curve solutions $w_1,w_2$,
Stokes extends by \eqref{eq:tracebound} to their $H^1$ difference and gives
\[
 E(w_1-w_2,w_1-w_2)
 =-\int_Y(w_1-w_2)(e^{\lambda w_1}-e^{\lambda w_2})\,d\mu\leq0.
\]
Their difference is constant; their unit exponential integrals force
that constant to be zero. This proves bounded curve uniqueness.

\section{Uniform selection and its order meaning}
The fixed bounded continuous potentials
\[
 L=Q-\sup Q\leq0,\qquad R=Q-\inf Q\geq0
\]
have curve measure $\mu$. They are respectively an exponential
subsolution and supersolution. For $z=(L-w_\lambda)_+$, the
$H^1$ chain rule and Stokes give
\[
 E(z,z)=-\int_Yz(1-e^{\lambda w_\lambda})\,d\mu\leq0,
\]
since $z>0$ implies $w_\lambda<L\leq0$.
To justify the test, the bounded density $e^{\lambda w_\lambda}$ and
\eqref{eq:tracebound} make the distributional right side continuous
on $H^1$. Smooth approximation extends Stokes to $z$; positive-part
chain rules identify its energy and canonical trace. Thus $z$ is
constant. A positive constant would give $w_\lambda<L$ everywhere and
$\int e^{\lambda w_\lambda}\,d\mu<1$, a contradiction. The identical
argument with $(w_\lambda-R)_+$ proves the upper bound. Consequently
\begin{equation}\label{eq:uniformbound}
 L\leq w_\lambda\leq R,\qquad
 \|w_\lambda\|_\infty\leq\operatorname{osc}Q.
\end{equation}

As $\lambda\downarrow0$, the actual equation gives
$\rho_\lambda=e^{\lambda w_\lambda}=1+O(\lambda)$ uniformly.
The Green identity and \eqref{eq:kernelbound} therefore imply
\[
 \|q_\lambda-Q\|_\infty
 \leq\|\rho_\lambda-1\|_\infty
       \sup_x\int_Y|G(x,y)|\,d\mu(y)=O(\lambda).
\]
The $q_\lambda$ are uniformly bounded, so
\[
 \lambda^{-1}\log Z_\lambda
  =\int_Yq_\lambda\,d\mu+O(\lambda)
  =\int_YQ\,d\mu+O(\lambda).
\]
This proves uniform convergence of the entire centered family with
the claimed rate. Supremum normalization preserves uniform convergence.

A genuine curve critical subsolution has non-pluripolar mass at most
one. If its measure dominates $\mu$, its mass is one and its measure
is exactly $\mu$. In dimension one its full positive current has
mass one, while the difference from the non-pluripolar measure is a
positive measure on its pole set. The difference has mass zero, so
the full current is $\mu$. Harmonic uniqueness makes the potential
$Q+C$. Its $\mu$-mean is nonpositive precisely for
$C\leq-\int Q\,d\mu$, giving the order characterization.

\section{Greatest exponential subsolutions and the measure threshold}
We use the bounded maximum principle of \cite[Lemma 1.2]{GL2} and
the strict domination principle of \cite[Lemma 4.2]{BGL}, whose Section 4
does not assume BMP. Their application here is to bounded potentials
on the strictly positive background $H$, with a constant strict barrier.
The inequalities compare actual Borel measures and require no density
against smooth volume. They do not supply existence for a singular measure;
the preceding variational argument supplies that separately.

Write $s=U_\lambda$. If $b$ is a bounded $H$-psh exponential subsolution,
then on $\{s<b-\epsilon\}$,
\[
 1_{\{s<b-\epsilon\}}\MA_H(s)
 \leq e^{-\lambda\epsilon}
       1_{\{s<b-\epsilon\}}\MA_H(b-\epsilon).
\]
Choose a constant $\rho\leq0$ below both bounded potentials, so that
$H+\ddc\rho=H$ dominates a Hermitian reference. Their difference is bounded,
and the displayed factor is strictly less than one. BGL Lemma 4.2 gives
$s\geq b-\epsilon$ everywhere, hence $s\geq b$. Applying the argument in
both directions proves uniqueness of bounded product solutions.

Now let $r$ be any genuine exponential subsolution. For $M>0$ put
$p=s-M$ and $W=\max(r,p)$. The finite supremum of $r$ and the bounded
lower floor $p$ make $W$ bounded. Choose $j$ with $-j<\inf p$, so
$r_j=\max(r,-j)$ is bounded $H$-psh and $W=\max(r_j,p)$. The bounded
maximum principle, including its equality-set term, gives
\[
 \MA_H(W)\geq
 1_{\{r_j>p\}}\MA_H(r_j)+1_{\{r_j\leq p\}}\MA_H(p).
\]
On the first set $r_j=r>-j$, and the actual strict-superlevel definition
identifies its measure with $\MA_H^{\rm ext}(r)$. On the second set,
$W=p$ and $\MA_H(p)=e^{\lambda s}\nu\geq e^{\lambda p}\nu$.
Thus $\MA_H(W)\geq e^{\lambda W}\nu$, and bounded comparison gives
$r\leq W\leq s$, canonically at every pole. Every product used before
restriction is bounded. No global unbounded-product or mass-invariance
theorem is used, even if the competitor has infinite mass.

Suppose $a>0$ and $\MA_H^{\rm ext}(v)\geq a\nu$ for a genuine $v$.
Set $B=\sup_Xv<\infty$. For each $\lambda>0$ the shifted potential
$v+\lambda^{-1}\log a-B$ is an exponential subsolution, since its
exponential is at most $a$. Greatestness therefore gives
\[
 v+\lambda^{-1}\log a-B
 \leq\lambda^{-1}\log c+w_\lambda.
\]
Evaluate at a fixed finite point of $v$, multiply by $\lambda$, and
use \eqref{eq:uniformbound} as $\lambda\downarrow0$. This proves $a\leq c$.
The actual finite equation for $u$ attains $c$, in the entire class, the
bounded ordinary type and the exact stretching class. This measure threshold
uses no bounded-density truncation. Greatestness alone does not prove
unbounded exponential uniqueness, and the critical order characterization
above remains confined to the separated sector.

\section{Actual product identities and singular examples}
The local bounded product recursion \cite[equation (2.1) and Corollary 2.2]{BT}, with a
common smooth correction for the nonclosed background
\cite[Section 2.1]{BGL}, gives exactly, for each bounded constructed curve $p$,
\[
 \MA_H^{\rm ext}(p)=N h_0^m\otimes(\kappa+\ddc p).
\]
Only one curve factor survives the dimension expansion. These are
actual Borel products; their equations and finite masses do not follow
merely from scalar convergence. Substitution of $Q$ and $w_\lambda$
proves all displayed equations, retaining the exponential supremum.
For every bounded $H$-psh $p$ and $t\geq1$, the genuine constant
$a_t=t\inf_Xp$ satisfies $a_t/t\leq p$. This proves the actual
stretching definition of $\Eenv(X,H,0)$.

Pull back bounded $h_0$-psh functions with masses tending to zero.
Their product masses are $N$ times those masses, proving $v_-(H)=0$.
A closed positive $h_0$ would have constant bounded mass $V_0$ by
Stokes, contradicting that conclusion; hence $h_0$ and $H$ are nonclosed.

A useful measure-only sufficient test for \eqref{eq:criterion} is
\[
 \mu(B_{\rm ch}(x,r))\leq C r^\delta
 \quad(0<r\leq1),\qquad \delta>0.
\]
Layer-cake bounds $\int_{K>L}K\,d\mu$ uniformly by
$C'(L+1)e^{-\delta L/2}$. Truncated kernel integrals are continuous;
their uniform convergence proves the criterion. In particular, on a
round circle take
\[
 d\mu(\vartheta)=(1+\epsilon\cos\vartheta)
                   \frac{d\vartheta}{2\pi},\qquad0<\epsilon<1.
\]
It has the bound with $\delta=1$ and is singular against smooth curve
volume; $\nu$ is singular against ambient smooth volume. Its logarithmic
potential is nonconstant on the circle, as the nonzero first Fourier
coefficient of the logarithmic kernel shows. Hence the constant-on-support
formula for unweighted circles does not supply these exponential solutions.
The variational construction supplies them.

The elliptic-quotient threefold factor with $V_0=24\pi$ and $v_-(h_0)=0$
has a complete construction and mass-collapse proof in \cite{Product}.
For this application fix its auxiliary disk probability as
$p(s)=Z^{-1}\exp(-1/(1-|s|^2))$ for $|s|<1$, extended by zero,
where $Z$ normalizes $p$ against $i\,ds\wedge d\bar s/(2\pi)$.
Then $\min_{|s|\leq1/4}p(s)=Z^{-1}e^{-16/15}>0$; this explicit
choice supplies the uniform positive block bound in that factor proof.
It makes $X$ a concrete fourfold and $c=96\pi/M$. The criterion concerns
this product geometry and these measures; arbitrary non-pluripolar
measures outside this regime remain open.

\begin{thebibliography}{9}
\bibitem{BT} E. Bedford and B. A. Taylor,
\emph{A new capacity for plurisubharmonic functions},
Acta Math. 149 (1982), 1--40.
\bibitem{BGL} S. Boucksom, V. Guedj and C. H. Lu,
\emph{Volumes of Bott--Chern classes}, arXiv:2406.01090v1.
\bibitem{GL2} V. Guedj and C. H. Lu,
\emph{Quasi-plurisubharmonic envelopes 2: Bounds on Monge--Amp\`ere volumes},
arXiv:2106.04272, Lemma 1.2.
\bibitem{BEGZ} S. Boucksom, P. Eyssidieux, V. Guedj and A. Zeriahi,
\emph{Monge--Amp\`ere equations in big cohomology classes},
Acta Math. 205 (2010), 199--262; arXiv:0812.3674, Section 1.2.
\bibitem{Dem} J.-P. Demailly,
\emph{Complex Analytic and Differential Geometry}, June 2012,
Chapter I, \S4.
\bibitem{Product} \emph{Exact equations for $L^1$ densities on products
with zero lower mass}, companion construction, Appendix A;
source: \texttt{P07-zero-lower-mass-product-equations.tex}.
\end{thebibliography}
\end{document}
